% !TEX root = ../main.tex
\lecture{7}{2026-11-02}{5.2 Ende}

\subsubsection{Zwischenwertsatz}
\begin{theorem}
    [Zwischenwertsatz] wenn eine Funktion [a,b] -> Z stetig, muss sie jeden Wert zwischen f(a) und f(b) annehmen
\end{theorem}

\subsubsection{Nullsatz}

\begin{theorem}
    \label{th_nullsatz}
    [Nullsatz] wenn Interval [a,b] worin f(a) < 0 und f(b) > 0, und die Funktion stetig dann muss es mindestens eine Nullstelle geben
\end{theorem}

